\newpage
\section{Identificación de tecnologías de desarrollo móvil}

Elegir tecnología me parece la decisión más importante de un proyecto y también la más cara de revertir: cambiar de framework con la app a medio hacer es casi empezar otra vez. Por eso en este apartado no me limito a listar opciones, sino que intento explicar por qué existe cada una y en qué situación concreta sería la mejor elección.

He mirado seis: las que son solución de facto en Android (Java y Kotlin, con Jetpack Compose como contexto) y las multiplataformas o nativas de iOS que más se usan (Flutter, React Native, .NET MAUI y Swift). El enunciado pide al menos cuatro, así que me ha dado para comparar de verdad en vez de listar sin criterio.

Antes de entrar, una aclaración que creo que evita bastante confusión: \textbf{Java y Kotlin no son competidoras la una de la otra}, porque en la práctica se usan juntas. Java es el lenguaje histórico de Android y hay millones de líneas de código ahí; Kotlin es el que Google recomienda para todo lo nuevo y está pensado justamente para convivir con el Java antiguo. Las alternativas reales entre sí son Kotlin, Flutter, React Native y MAUI.

\subsection{Java (desarrollo nativo de Android)}

\campo{Nombre de la tecnología.} Java, el lenguaje original y oficial del desarrollo Android.

\campo{Lenguaje.} Java (y Kotlin, que convive con él).

\campo{Tipo de desarrollo.} Nativo. El código compila a bytecode y se ejecuta en la máquina virtual de Android (ART), sobre las APIs del sistema.

\campo{Sistemas operativos compatibles.} Android. Y aquí está lo que hay que tener claro: solo Android. Si compilas con Java para Android no te sirve para nada en iOS.

\campo{Principales características.}
\begin{itemize}[leftmargin=*,itemsep=0.2em,topsep=0.2em]
  \item Acceso completo y directo a todas las APIs del sistema (cámara, GPS, Bluetooth, sensores, notificaciones, Wear OS\ldots).
  \item Interfaz de dos formas: con \textbf{vistas XML} (\code{res/layout/*.xml}) o de forma declarativa con \textbf{Jetpack Compose}.
  \item Herramienta oficial: \textbf{Android Studio}, que trae compilador (Gradle), emulador, inspector de memoria y de red, y un perfilador de CPU.
  \item Un dato que me llamó la atención: en mayo de 2026 Google publicó una entrada de blog titulada literalmente \enquote{Android UI development is Compose-first}. El sistema de vistas clásico entra en modo mantenimiento.
\end{itemize}

\campo{Ventajas.} Rendimiento máximo, porque no hay capas intermedias; cobertura total de la plataforma; ecosistema maduro; estabilidad a largo plazo.

\campo{Desventajas.} Solo sirve para Android; mucho código repetitivo con el sistema de vistas; fragmentación, que son miles de fabricantes; mantenimiento costoso.

\campo{Ejemplo o situación adecuada.} Una app de banca o de salud que necesite leer el NFC de la tarjeta sanitaria, abrir la cámara para escanear recetas y funcionar en el \emph{longtail} de Android barato, donde el rendimiento y el tamaño del APK importan mucho.

\subsection{Kotlin (desarrollo nativo de Android, moderno)}

\campo{Nombre de la tecnología.} Kotlin, el lenguaje oficial preferido por Google para Android.

\campo{Lenguaje.} Kotlin.

\campo{Tipo de desarrollo.} Nativo para Android y además multiplataforma mediante \textbf{Kotlin Multiplatform (KMP)}.

\campo{Sistemas operativos compatibles.} Android de forma nativa; y mediante KMP, la misma lógica de negocio se puede compartir con iOS, macOS, Windows, Linux, web (Wasm/JS), tvOS y watchOS.

\campo{Principales características.}
\begin{itemize}[leftmargin=*,itemsep=0.2em,topsep=0.2em]
  \item Conciso y seguro: el tipado nulo es del lenguaje, y eso ya elimina muchísimos errores en tiempo de ejecución.
  \item Coroutines, que es la forma oficial de hacer asincronía sin bloquear la interfaz.
  \item Interoperabilidad total con Java: puedes llamar a Java desde Kotlin y al revés.
  \item Jetpack Compose para la interfaz moderna declarativa. Y también Compose Multiplatform.
  \item Google mantiene una estrategia \textbf{Kotlin-first}: la documentación nueva sale primero en Kotlin.
\end{itemize}

\campo{Ventajas.} Menos código y menos errores que Java; las coroutines resuelven muy bien el trabajo en segundo plano; KMP permite reutilizar la lógica entre Android e iOS; herramientas excelentes.

\campo{Desventajas.} Curva de aprendizaje; KMP todavía está madurando; interoperabilidad con el ecosistema antiguo; dependencia de Google.

\campo{Ejemplo o situación adecuada.} Una app de comercio electrónico que nace para Android, tiene una lógica de precios, stock y carrito complicada, y en un año quiere publicar en iOS. Con KMP se comparte el 70-80\,\% del código de negocio.

\subsection{Flutter}

\campo{Nombre de la tecnología.} Flutter, el framework multiplataforma de Google.

\campo{Lenguaje.} Dart.

\campo{Tipo de desarrollo.} Multiplataforma, con una única base de código.

\campo{Sistemas operativos compatibles.} Android (API 24+), iOS/iPadOS (15.0+), Windows, macOS, Linux y web. Es la lista más amplia de las seis.

\campo{Principales características.}
\begin{itemize}[leftmargin=*,itemsep=0.2em,topsep=0.2em]
  \item Flutter no usa los botones ni las vistas nativas: dibuja todos los píxeles él mismo a través de su propio árbol de widgets.
  \item Impeller es su renderizador por defecto desde la 3.27 en iOS, y en Android a partir de la API 29.
  \item Hot reload de verdad: cambias un \code{Text} y lo ves en el móvil en menos de un segundo.
  \item Ecosistema en \code{pub.dev} con más de 1.300 millones de descargas en los últimos 30 días.
  \item Versión actual: Flutter 3.47 (agosto de 2026).
\end{itemize}

\campo{Ventajas.} Una sola base de código para todo; interfaz visual idéntica en todas partes; rendimiento muy cercano al nativo; hot reload excelente; consistencia.

\campo{Desventajas.} La aplicación resultante pesa más; dependencia de plugins; el aspecto nativo no es exacto; Dart es un lenguaje \enquote{nuevo}; riesgo de quedarte atado a una sola plataforma.

\campo{Ejemplo o situación adecuada.} Una startup que quiere lanzar la misma app en iOS, Android, web y escritorio en tres meses con un equipo pequeño.

\subsection{React Native}

\campo{Nombre de la tecnología.} React Native, el framework multiplataforma de Meta.

\campo{Lenguaje.} JavaScript o TypeScript.

\campo{Tipo de desarrollo.} Multiplataforma, sobre todo Android e iOS.

\campo{Sistemas operativos compatibles.} Android, iOS y, mediante proyectos externos (\emph{out-of-tree}), macOS, Windows, visionOS, tvOS, web y HarmonyOS.

\campo{Principales características.}
\begin{itemize}[leftmargin=*,itemsep=0.2em,topsep=0.2em]
  \item Al revés que Flutter, sí usa las vistas nativas de cada sistema.
  \item La Nueva Arquitectura (Fabric y Turbo Native Modules) es el modo por defecto desde la 0.76, y desde la 0.82 (octubre de 2025) es la única arquitectura.
  \item Fast Refresh activado por defecto.
  \item Versión actual: React Native 0.87 (agosto de 2026), que exige Node 22+.
  \item Ecosistema npm: cualquier librería de JavaScript probablemente ya tiene bindings.
\end{itemize}

\campo{Ventajas.} Enorme mercado de desarrolladores; ecosistema npm gigantesco; integración incremental; buen rendimiento con la Nueva Arquitectura; plataformas adicionales.

\campo{Desventajas.} Los módulos nativos son fuente de problemas; las actualizaciones mayores duelen; rendimiento variable; depurar es más complejo; tiene menos presencia que Flutter en el mercado de frameworks.

\campo{Ejemplo o situación adecuada.} Una empresa que ya tiene un equipo web grande y quiere sacar la app móvil sin contratar gente nueva.

\subsection{.NET MAUI}

\campo{Nombre de la tecnología.} .NET MAUI (Multi-platform App UI), el sucesor de Xamarin.Forms dentro del ecosistema .NET de Microsoft.

\campo{Lenguaje.} C\#.

\campo{Tipo de desarrollo.} Multiplataforma.

\campo{Sistemas operativos compatibles.} Android (API 24+), iOS 11+, macOS (Mac Catalyst), Windows 10 1809+ con WinUI 3, y tvOS.

\campo{Principales características.}
\begin{itemize}[leftmargin=*,itemsep=0.2em,topsep=0.2em]
  \item Usa los controles nativos de cada plataforma mediante un modelo de handlers.
  \item XAML para declarar la interfaz y C\# para la lógica.
  \item Hot Reload, incluido el XAML Hot Reload.
  \item Estado actual: .NET MAUI 10 (noviembre de 2025). .NET MAUI 9 dejó de soportarse en mayo de 2026 y la versión 10 termina su soporte en mayo de 2027.
\end{itemize}

\campo{Ventajas.} Un solo lenguaje y un solo IDE (Visual Studio); controles nativos reales; buena integración con .NET; cubre escritorio; C\# es muy productivo.

\campo{Desventajas.} Las críticas de la comunidad sobre estabilidad y fiabilidad son serias y están documentadas; ecosistema de UI más pequeño; rendimiento irregular en apps grandes; ciclo de soporte más corto (12 meses por versión); comunidad más pequeña.

\campo{Ejemplo o situación adecuada.} Una empresa mediana con un equipo de C\#/.NET que ya tiene una aplicación de escritorio Windows y quiere sacar la versión móvil y la de macOS con el mismo equipo.

\subsection{Swift (desarrollo nativo de iOS)}

\campo{Nombre de la tecnología.} Swift, el lenguaje oficial de Apple.

\campo{Lenguaje.} Swift.

\campo{Tipo de desarrollo.} Nativo, y solo del ecosistema Apple.

\campo{Sistemas operativos compatibles.} iOS, iPadOS, macOS, tvOS, visionOS y watchOS. No hay versión para Android.

\campo{Principales características.}
\begin{itemize}[leftmargin=*,itemsep=0.2em,topsep=0.2em]
  \item Dos formas de hacer interfaz: SwiftUI (declarativo) y UIKit (imperativo). Apple dice que conviven y que se adopte \enquote{a tu propio ritmo}.
  \item Concurrencia moderna: la versión actual es Swift 6.4 (septiembre de 2026), con \emph{strict concurrency}.
  \item Xcode Previews: ves la pantalla en el canvas lateral sin arrancar la app.
  \item Integración profunda: Core Data, SwiftUI, Metal, Core ML, HealthKit, ARKit\ldots
  \item Seguridad de memoria gestionada (ARC).
\end{itemize}

\campo{Ventajas.} Máximo rendimiento y fidelidad en plataforma Apple; acceso completo a cada API; lenguaje moderno y seguro; herramientas excelentes (Xcode + SwiftUI Previews); soporte a largo plazo coherente.

\campo{Desventajas.} Lock-in total; cuesta aprender dos cosas (SwiftUI y UIKit); ecosistema de herramientas más estrecho; la API evoluciona muy rápido.

\campo{Ejemplo o situación adecuada.} Cualquier app de calidad o de prestigio en el ecosistema Apple. Siri, Things, Fantastical, Bear, Overcast\ldots todas están escritas en Swift.